\section{Fundamentação teórica}

As Large Language Models (LLMs) são modelos de inteligência artificial (IA) treinados sobre grandes volumes de dados textuais com o objetivo de compreender e gerar linguagem natural \cite{born2025revisao}. Com base na arquitetura Transformer, esses modelos foram capazes de aprender e se especializar por meio de padrões linguísticos e relações semânticas que lhes permitem realizar tarefas como geração de texto, tradução, respostas a perguntas complexas e produção de código. Alguns exemplos utilizados atualmente incluem o GPT, Claude, Gemini e Llama \cite{minaee_large_2024}.

No meio da Engenharia de Software (ES), esses modelos foram introduzidos inicialmente como assistentes para geração de pequenos trechos de código, passando a apoiar diversas atividades ao longo do ciclo de desenvolvimento. Na Engenharia de Requisitos (ER), por sua vez, auxiliam na elicitação, refinamento e validação dos requisitos \cite{pires_prompts_2025}.

Com o avanço da tecnologia, as LLMs começaram a atuar de forma mais autônoma, diretamente do terminal do computador de um desenvolvedor. Desse modo, assumiram papéis especializados para execução das tarefas, colaborando na realização de atividades complexas de desenvolvimento de software \cite{ribeiro_software_2025}.

Atualmente, o principal desafio enfrentado ao integrar as LLMs no processo de codificação são as alucinações, sendo uma limitação crítica dos modelos. Como forma de contornar o problema, pode-se usar de estratégias como Retrieval Augmented Generation (RAG), configuração de modelos mais robustos, feedback contínuo humanizado e prompts mais específicos \cite{minaee_large_2024}.

Dentre os benefícios observados na literatura, observa-se o aumento da produtividade dos desenvolvedores, redução do tempo necessário para execução de tarefas repetitivas, apoio na documentação, auxílio na geração de testes e identificação de problemas e democratização do acesso ao conhecimento técnico \cite{tolio_llms_2026}.

Diante desse cenário de adoção crescente, a Engenharia de Prompt (EP) vem se tornando uma ferramenta primordial para analisar e mitigar as alucinações geradas pelos modelos de inteligência artificial \cite{silva_engenharia_2024}. Ela contribui na criação de um texto elaborado, propondo instruções estruturadas para orientar as LLMs na execução de tarefas específicas. A qualidade e a clareza no prompt estão diretamente ligadas à consistência e à relevância das respostas geradas.

Diversas técnicas têm sido propostas para melhorar o desempenho das respostas \cite{nacimento2024exploraccao}. Entre elas destacam-se o \textit{Role Prompting}, no qual o modelo recebe um papel específico a ser desempenhado, como analista de requisitos, PO ou arquiteto de software. O \textit{Chain-of-Thought}, que incentiva a geração de um raciocínio intermediário para analisar e pensar como será feita a resposta final \cite{minaee_large_2024}. O \textit{Structured Prompting}, baseado na definição de um formato estruturado de entrada e saída, reduzindo ambiguidade e tornando as respostas mais consistentes \cite{ciudrowski2026explorando}. Por último, o \textit{Few-Shot Prompting}, que se utiliza de exemplos para demonstrar o comportamento esperado do modelo \cite{minaee_large_2024}.

Essas técnicas de prompt costumam ser combinadas com estratégias mais amplas de organização das LLMs, como os sistemas encadeados, chamados na literatura de chaining, que consistem na utilização de vários modelos especializados realizando atividades conjuntas, de forma que atinjam um objetivo comum \cite{minaee_large_2024}.

Diferente da estratégia de multiagente dinâmico, o chaining não envolve diálogo e refinamento iterativo entre as IAs. Ele opera de forma linear, com auxílio de um humano para realizar a passagem de conhecimento de um elo para o outro \cite{minaee_large_2024}.

Na Engenharia de Software, essa prática tem se tornado cada vez mais utilizada para distribuir responsabilidades entre especialistas de IA. Sua abordagem busca produzir a divisão de tarefas em partes menores, na busca de resultados mais elaborados e assertivos \cite{larguesa2025engenharia}. Sua utilização em conjunto com LLMs especializadas para refinamento do prompt e codificação pode potencializar os resultados \cite{ribeiro_software_2025}. Desse modo, é possível produzir um artefato que serve como entrada para agentes responsáveis pela implementação do software.
